\section*{第 \ref{ch:b1:replication-mitosis} 章——DNA 复制与有丝分裂}

\begin{solution}{exo:b1:replication-mitosis:1}
聚合酶只能将\omterm{def:b1:nucleic-acids:nucleotide}{核苷酸}加到 $3'$ 端，因此新链
沿 $5' \to 3'$ 方向延伸；两条模板反向平行，故在\omterm{def:b1:replication-mitosis:fork}{复制叉}处，一条新链
可朝\omterm{def:b1:replication-mitosis:fork}{复制叉}方向连续延伸，另一条则须背离\omterm{def:b1:replication-mitosis:fork}{复制叉}延伸，
随着更多模板暴露，
以片段形式重新起始。
\end{solution}

\begin{solution}{exo:b1:replication-mitosis:2}
解旋酶（解开双链）、单链结合蛋白（使两链保持分离）、
拓扑异构酶（松解前方扭转）、引发酶（合成 \omterm{def:b1:nucleic-acids:chain}{RNA} 引物）、
\omterm{def:b1:nucleic-acids:chain}{DNA} 聚合酶（延伸引物并校对）、第二种聚合酶
（替换引物）、连接酶（封闭片段）。
\end{solution}

\begin{solution}{exo:b1:replication-mitosis:3}
G$_1$ 期：$q$，2n 条\omterm{def:b1:genomes:genome}{染色体}，每条含一条\omterm{def:b1:replication-mitosis:mitosis}{染色单体}。G$_2$ 期：$2q$，2n
条\omterm{def:b1:genomes:genome}{染色体}，每条含两条\omterm{def:b1:replication-mitosis:mitosis}{染色单体}。\omterm{def:b1:replication-mitosis:mitosis}{有丝分裂}后：每个子\omterm{def:b1:cell-unit-of-life:cell}{细胞} $q$，2n
条\omterm{def:b1:genomes:genome}{染色体}。
\end{solution}

\begin{solution}{exo:b1:replication-mitosis:4}
前期（\omterm{def:b1:genomes:genome}{染色体}凝缩，\omterm{def:b1:replication-mitosis:mitosis}{纺锤体}形成），前中期（\omterm{def:b1:eukaryotic-cell:nucleus}{核膜}破裂，
动粒附着），中期（\omterm{def:b1:genomes:genome}{染色体}排列在赤道板上），后期
（\omterm{def:b1:replication-mitosis:mitosis}{染色单体}分离），末期（\omterm{def:b1:eukaryotic-cell:nucleus}{核膜}重新形成），随后
胞质分裂。
\end{solution}

\begin{solution}{exo:b1:replication-mitosis:5}
\omterm{def:b1:genomes:genome}{基因组}的一半由后随链合成：$\num{6.4e9}/150 =
\num{4.3e7}$ 个片段；引物、引物去除与连接的次数均与之相同。
\end{solution}

\begin{solution}{exo:b1:replication-mitosis:6}
$\num{6.4e9}\times 10^{-5} = \num{64000}$；$\times 10^{-7}$：640；
$\times 10^{-9}$：每次分裂约 6 个突变。
\end{solution}

\begin{solution}{exo:b1:replication-mitosis:7}
$5000/80 = 62$ 次分裂。肿瘤须能无限分裂；若无
\omterm{def:b1:genomes:karyotype}{端粒}酶，其\omterm{def:b1:genomes:karyotype}{端粒}将耗尽，\omterm{def:b1:cell-unit-of-life:cell}{细胞}将停滞或
死亡，因此几乎所有癌症都会重新激活该酶。
\end{solution}

\begin{solution}{exo:b1:replication-mitosis:8}
\omterm{def:b1:replication-mitosis:cycle}{细胞周期} $= 1\,\mathrm{h}/0.05 = \qty{20}{h}$；S 期 $= 0.30\times 20 =
\qty{6}{h}$。
\end{solution}

\begin{solution}{exo:b1:replication-mitosis:9}
G$_1$ 期：8 条\omterm{def:b1:genomes:genome}{染色体}，8 条\omterm{def:b1:replication-mitosis:mitosis}{染色单体}，$q$。G$_2$ 期：8、16、$2q$。中期：
8、16、$2q$。子\omterm{def:b1:cell-unit-of-life:cell}{细胞}：8、8、$q$。
\end{solution}

\begin{solution}{exo:b1:replication-mitosis:10}
每隔 \qty{20}{min} 即有新一轮复制在复制起点启动，此时前一轮
尚未结束，复制相互重叠，新生\omterm{def:b1:cell-unit-of-life:cell}{细胞}的 \omterm{def:b1:nucleic-acids:chain}{DNA} 已部分复制。
分裂时，即将完成复制的\omterm{def:b1:genomes:genome}{染色体}，其复制启动于
\qty{40}{min} 之前；第二轮复制启动于 \qty{20}{min} 之前，
\omterm{def:b1:replication-mitosis:fork}{复制叉}已推进到一半；第三轮刚刚开始，
因此分裂\omterm{def:b1:cell-unit-of-life:cell}{细胞}中有 4 个复制起点、6 个\omterm{def:b1:replication-mitosis:fork}{复制叉}。
\end{solution}

\begin{solution}{exo:b1:replication-mitosis:11}
不形成\omterm{def:b1:replication-mitosis:mitosis}{纺锤体}：\omterm{def:b1:genomes:genome}{染色体}虽凝缩，却无法排列或
分离，\omterm{def:b1:cell-unit-of-life:cell}{细胞}会在\omterm{def:b1:replication-mitosis:mitosis}{纺锤体}检验点停滞，呈
类似中期的状态。若\omterm{def:b1:cell-unit-of-life:cell}{细胞}未分裂便回到\omterm{def:b1:replication-mitosis:cycle}{间期}，
则其\omterm{def:b1:genomes:genome}{染色体}数为 4n（四倍体）。停滞的\omterm{def:b1:cell-unit-of-life:cell}{细胞}具有
凝缩且分开的\omterm{def:b1:genomes:genome}{染色体}，这正是\omterm{def:b1:genomes:karyotype}{核型}分析所需要的，
故施用该药物以使其积累。
\end{solution}

\begin{solution}{exo:b1:replication-mitosis:12}
复制保证每条\omterm{def:b1:genomes:genome}{染色体}形成两条完全相同的\omterm{def:b1:replication-mitosis:mitosis}{染色单体}，
误差率仅为十亿分之一；\omterm{def:b1:replication-mitosis:mitosis}{纺锤体}保证每个子\omterm{def:b1:cell-unit-of-life:cell}{细胞}各获得每对中的一条，
其方式是将每条\omterm{def:b1:genomes:genome}{染色体}与两极相连，
并在所有\omterm{def:b1:genomes:genome}{染色体}均受张力之前阻止后期开始。
G$_2$ 期检验点确保在开始计数之前复制已完成；
\omterm{def:b1:replication-mitosis:mitosis}{纺锤体}检验点确保在分离之前计数已就绪。
分离出错的\omterm{def:b1:cell-unit-of-life:cell}{细胞}，\omterm{def:b1:genomes:genome}{染色体}会多一条或少一条（非整倍性）：
拷贝准确无误，却被送到了错误的地址——
这类\omterm{def:b1:cell-unit-of-life:cell}{细胞}多数会死亡，而非整倍性
恰是多数癌症的特征。
\end{solution}

\begin{solution}{pb:b1:replication-mitosis:1}
\textbf{1.} $50\times 28\,800 = \qty{1.44e6}{bp}$。
\textbf{2.} \qty{2.88e6}{bp}。
\textbf{3.} $\num{6.4e9}/\num{2.88e6} = 2220$ 个复制起点。
\textbf{4.} \qty{2.9}{Mb}，即 $2.9\times 10^6\times 0.34\,\mathrm{nm} =
\qty{0.98}{mm}$ 的 \omterm{def:b1:nucleic-acids:chain}{DNA} 位于两复制起点之间。
\textbf{5.} $\num{6.4e9}/\num{40000} = \qty{160}{kb}$ 间距；每个
\omterm{def:b1:replication-mitosis:fork}{复制叉}约移动 \qty{80}{kb}，约相当于半小时的工作量。
\textbf{6.} 以 \qty{50}{bp/s} 的速率复制 \qty{125}{Mb}，需 $\num{2.5e6}$ 秒，
29 天。
\textbf{7.} \omterm{def:b1:replication-mitosis:fork}{复制叉}速度受化学过程及
\omterm{def:b1:genomes:nucleosome}{染色质}的限制（相较细菌慢了二十倍）；复制起点可以无限增加，
因此缩短 S 期应靠增加起始点，
而非加快\omterm{def:b1:replication-mitosis:fork}{复制叉}。
\textbf{8.} 后随链合成覆盖一半\omterm{def:b1:genomes:genome}{基因组}，即 $\num{3.2e9}$
个\omterm{def:b1:nucleic-acids:nucleotide}{核苷酸}，形成 $\num{2.1e7}$ 个片段（每个复制起点的两个\omterm{def:b1:replication-mitosis:fork}{复制叉}各有
一条后随链，两半合起来恰为一个\omterm{def:b1:genomes:genome}{基因组}的量）。
\textbf{9.} $\num{2.1e7} + 2\times\num{40000} = \num{2.1e7}$ 个引物。
\textbf{10.} 约 $\num{2.1e7}$ 次连接（每个片段一次），另加
每次\omterm{def:b1:replication-mitosis:fork}{复制叉}相遇一次。
\textbf{11.} $2\times\num{6.4e9} = \num{1.28e10}$ 个\omterm{def:b1:nucleic-acids:nucleotide}{核苷酸}。
\textbf{12.} $\num{2.56e10}$ 个 \omterm{def:b1:water-small-molecules:atp}{ATP}；在 \qty{28800}{s} 内，$\num{8.9e5}$
每秒。
\textbf{13.} 约占\omterm{def:b1:cell-unit-of-life:cell}{细胞} \omterm{def:b1:water-small-molecules:atp}{ATP} 周转量的 \qty{0.1}{\%}：
聚合反应本身代价低廉；昂贵的是合成
\omterm{def:b1:nucleic-acids:nucleotide}{核苷酸}和\omterm{def:b1:genomes:nucleosome}{组蛋白}。
\textbf{14.} \num{64000}、640 和 6.4 个误差。
\textbf{15.} $1500\times 10^{-9} = \num{1.5e-6}$：每次分裂中，约每七
十万个\omterm{def:b1:genomes:gene}{基因}中有一个发生突变。
\textbf{16.} $10^{16}\times 6.4 = \num{6.4e16}$ 个突变——\omterm{def:b1:genomes:genome}{基因组}中每一种
可能的单点改变均多次重复出现，散布于
$10^{13}$ 个\omterm{def:b1:cell-unit-of-life:cell}{细胞}中：几乎所有突变都落在非编码 \omterm{def:b1:nucleic-acids:chain}{DNA} 上，或发生在
将被脱落的\omterm{def:b1:cell-unit-of-life:cell}{细胞}中，突变\omterm{def:b1:cell-unit-of-life:cell}{细胞}仅是数十亿个之一；只有少数
突变组合在同一\omterm{def:b1:cell-unit-of-life:cell}{细胞}中同时出现时才有意义。
\textbf{17.} 增至百倍（$10^{-7}$）：导致癌症的突变序列，
在 $10^{-9}$ 下需数十年才能积累，
则积累快得多——遗传的错配修复缺陷
会导致成年早期的结肠癌。
\textbf{18.} 每个\omterm{def:b1:replication-mitosis:fork}{复制叉}复制 $\num{2.3e6}$ 个\omterm{thm:b1:nucleic-acids:helix}{碱基对}，速率为 \qty{1000}{bp/s}：
\qty{2300}{s}，即 \qty{38}{min}。
\textbf{19.} 约有两轮复制相互重叠；新生\omterm{def:b1:cell-unit-of-life:cell}{细胞}带有
两个复制起点（各已复制一次），\omterm{def:b1:replication-mitosis:fork}{复制叉}已推进到一半：
其\omterm{def:b1:genomes:genome}{染色体}出生时已部分复制。
\textbf{20.} $\num{4.6e6}/1500 = 3070$ 个片段（每轮复制），约
每秒 1.3 个。
\textbf{21.} $\num{4.6e6}\times 10^{-9} = \num{4.6e-3}$ 个误差，
每轮复制：约每 200 个子\omterm{def:b1:cell-unit-of-life:cell}{细胞}中有一个携带新突变。
\textbf{22.} $10^9\times\num{4.6e-3} = \num{4.6e6}$ 个新突变，每
毫升；某一给定的 \qty{1}{kb} \omterm{def:b1:genomes:gene}{基因}平均受到 $\num{4.6e6}\times
1000/\num{4.6e6} = 1000$ 次。
\textbf{23.} 由于每毫升每次分裂中，任一给定\omterm{def:b1:genomes:gene}{基因}都有一千个独立突变，
故每一种可能的单点改变——包括
赋予抗性的那些——在施用抗生素之前便已存在于
大量培养物中；药物起的是选择作用，而非
产生突变。
\textbf{24.} 人类：$\num{6.4e9}$ 个\omterm{thm:b1:nucleic-acids:helix}{碱基对}，\num{40000} 个复制起点，
\qty{50}{bp/s}，\qty{8}{h}。细菌：$\num{4.6e6}$ 个\omterm{thm:b1:nucleic-acids:helix}{碱基对}，一个复制起点，
\qty{1000}{bp/s}，\qty{40}{min}。\omterm{def:b1:genomes:genome}{基因组}大一千倍，
\omterm{def:b1:replication-mitosis:fork}{复制叉}速度慢二十倍，所用时间却仅为十二倍：
差别在于数以万计的复制起点
并行工作。
\textbf{25.} 若全部同时启动，至少约需 \num{2200} 个复制起点；
实际约有 \num{40000} 个，每 \qty{160}{kb} 一个，依次启动，
使每个\omterm{def:b1:replication-mitosis:fork}{复制叉}约移动 \qty{80}{kb}。
\end{solution}
